\documentclass[11pt]{article}
\usepackage[T1]{fontenc}
\usepackage{lmodern}
\usepackage[margin=0.9in]{geometry}
\usepackage{amsmath,amssymb,amsthm}
\usepackage{microtype}
\usepackage[hidelinks]{hyperref}
\usepackage{enumitem}
\setlength{\parindent}{0pt}
\setlength{\parskip}{6pt}
\newtheorem{theorem}{Theorem}
\newcommand{\GL}{\mathrm{GL}}
\newcommand{\Fp}{\mathbb F_p}
\newcommand{\ord}{\operatorname{ord}}
\begin{document}
\begin{center}
{\Large\bfseries Disproof of conjecture 00000000161}\\[5pt]
{\large A uniform prime-divisor bound in dimension four}\\[7pt]
C0ldSmi1e \quad --- \quad AI-assisted submission\\
9 October 2026
\end{center}

\paragraph{Claim and scope.}
Both language versions assert that, for every fixed integer $n\ge2$, a uniformly
chosen $A\in\GL_n(\Fp)$ has, with probability tending to one as the prime $p$
tends to infinity, a prime divisor of $\ord(A)$ exceeding
$p^{n-1}/\operatorname{poly}(n)$. The conventional denominator is positive and
independent of $p$. No particular polynomial is specified. We disprove the
claim for every possible positive denominator value at $n=4$.

\begin{theorem}
Let $C>0$. For every prime $p$ satisfying $p\ge3C$ and every
$A\in\GL_4(\Fp)$, no prime divisor of $\ord(A)$ exceeds $p^3/C$.
Consequently the uniform probability of the conjectured event is eventually
exactly zero, and its limit along primes is zero, not one.
\end{theorem}

\begin{proof}
Counting ordered linearly independent columns gives
\[
 |\GL_4(\Fp)|=\prod_{i=0}^{3}(p^4-p^i)
   =p^6(p-1)^4(p+1)^2(p^2+1)(p^2+p+1).
\]
For completeness, the individual factors are
\begin{align*}
 p^4-1&=(p-1)(p+1)(p^2+1),&
 p^4-p&=p(p-1)(p^2+p+1),\\
 p^4-p^2&=p^2(p-1)(p+1),&
 p^4-p^3&=p^3(p-1).
\end{align*}
If a prime $r$ divides $\ord(A)$, Lagrange's theorem implies that it divides
$|\GL_4(\Fp)|$. Euclid's lemma therefore makes $r$ a divisor of at least one of
$p$, $p-1$, $p+1$, $p^2+1$, and $p^2+p+1$. These integers are positive because
$p\ge2$, and each is at most $p^2+p+1$. Thus
\[
 r\le p^2+p+1\le3p^2\le\frac{p^3}{C}\qquad(p\ge3C).
\]
The conjecture requires a \emph{strictly} larger prime divisor, so its event
is empty. Its uniform counting probability is zero. Primes are unbounded;
hence this holds eventually along the prime-indexed limit. Uniqueness of real
limits on a nontrivial filter rules out a limit of one.
\end{proof}

\paragraph{Why this refutes the full claim.}
For any proposed positive polynomial $P$, put $C=P(4)>0$. The conjecture includes
the dimension $4$, where $n-1=3$, and is contradicted by the theorem. The same
argument rules out even an arbitrary positive function of dimension as the
denominator. No conclusion about other dimensions, or about a denominator
allowed to depend on $p$, is required.

\newpage
\section*{Formal correspondence and reproduction}

The exact bilingual original is included as \texttt{conjecture.md}. The Lean
project uses the original threshold, the actual general linear group, and the
actual group-theoretic element order. The relevant definitions are:
\begin{itemize}[leftmargin=1.4em,itemsep=3pt,topsep=3pt]
\item \texttt{PrimeIndex} is $\{p\in\mathbb N:p\text{ is prime}\}$ with the
usual inherited order. Its \texttt{atTop} filter is nontrivial; coercion to
$\mathbb N$ tends to infinity by Euclid's theorem.
\item \texttt{GeneralLinear n p} is
\texttt{Matrix.GeneralLinearGroup (Fin n) (ZMod p)}. For prime $p$ this is
$\GL_n(\Fp)$, the units of the square matrix ring over the $p$-element field.
\item \texttt{LargePrimeFactor} requires a natural prime $r$ dividing
\texttt{orderOf A} and satisfying $p^{n-1}/C<r$.
\item \texttt{successProbability} is the cardinality of that event divided
by the cardinality of the full group. Positivity of the finite sample-space
cardinality and membership of this ratio in $[0,1]$ are proved.
\item \texttt{OriginalClaim} explicitly means that there exists a real
polynomial $P$, positive at every integer $n\ge2$, for which the corresponding
probability tends to one at every such fixed dimension.
\end{itemize}

The mathematical files are \texttt{Conjecture161.lean} and
\texttt{Conjecture161/Arithmetic.lean}. The key bridge,
\nolinkurl{prime_divisor_order_le}, applies the arithmetic factorization to
Mathlib's \nolinkurl{Matrix.card_GL_field} and
\nolinkurl{orderOf_dvd_natCard}. Thus the group cardinality and its link to
individual element orders are proved within Lean, not merely supplied as
numerical assumptions.

The final theorem is
\begin{center}
\texttt{Conjecture161.conjecture\_false :}\quad
\texttt{$\neg$ Conjecture161.OriginalClaim}.
\end{center}
The stronger results
\nolinkurl{successProbability_four_not_tendsto_one} and
\nolinkurl{no_positive_denominator_function} make the independence from the
choice of polynomial explicit. Exact event emptiness, eventual zero, and the
limit zero are all separately formalized.

\paragraph{Verification.}
The project pins Lean 4.19.0 and Mathlib v4.19.0. The Mathlib commit is
\begin{center}
\texttt{c44e0c8ee63ca166450922a373c7409c5d26b00b}.
\end{center}
All eight transitive dependency
revisions are pinned in \texttt{lake-manifest.json}. With the pinned Lean on
\texttt{PATH}, obtain dependencies using \texttt{lake exe cache get}, then run
\texttt{python3 verify.py} inside \texttt{lean/}.
The verifier checks the pins, rebuilds the project's own generated objects,
replays all Lean files with warnings treated as errors, and prints the full
types and axiom dependencies of all 14 named mathematical theorems.
The recorded author run passed. Its axiom sets contain only
\texttt{propext}, \texttt{Classical.choice}, and \texttt{Quot.sound}.
There are no admitted proofs, added axioms, or \texttt{native\_decide} calls.
No auxiliary numerical computation is needed: the proof is uniform in every
prime and every group element under the displayed bound.

\paragraph{Provenance and review status.}
The proof author worked from the original bilingual statement, contribution
rules, and stock pinned libraries. After independently deriving the argument,
the author delegated only the generic arithmetic factorization lemma to a
fresh helper; the full disclosure is in \texttt{AUTHOR\_NOTES.txt}. The package's
verification records distinguish author checks, independent internal review,
and repository eligibility checks. These records do not constitute maintainer
acceptance.
\end{document}
